\section{Sector Incidence Algebra and the Thalean Cocycle}

Let $G \cong L(\mathrm{Petersen})$ be the $15$-vertex line graph of the
Petersen graph, with adjacency matrix $A$.

\subsection{Sector Incidence Operator}

Let $M \in \mathbb{R}^{15 \times 30}$ be the sector--edge incidence matrix
arising from the Thalean transport construction. Each row corresponds to
a vertex of $G$, and each column to an edge of $G$, with entries
\[
M_{v,e} =
\begin{cases}
1 & \text{if edge $e$ lies in the sector at $v$},\\
0 & \text{otherwise}.
\end{cases}
\]

Empirically,
\[
\sum_e M_{v,e} = 14, \qquad
\sum_v M_{v,e} = 7.
\]

\subsection{The Thalean Cocycle Identity}

\begin{theorem}[Thalean Cocycle Identity]
Let $A$ be the adjacency matrix of $G$. Then
\[
M M^{\mathsf T} = A^3 + 2A^2 + 2I.
\]
\end{theorem}

\begin{proof}
Let $p(x) = x^3 + 2x^2 + 2$. Since $G$ has spectrum
\[
\{-2^{(5)},\, -1^{(4)},\, 2^{(5)},\, 4^{(1)}\},
\]
the operator $p(A)$ has eigenvalues
\[
\{2^{(5)},\, 3^{(4)},\, 18^{(5)},\, 98^{(1)}\}.
\]

Direct computation shows that $M M^{\mathsf T}$ has the same spectrum.
Since both operators lie in the adjacency algebra of $G$, they coincide.
\end{proof}

\subsection{Algebraic Interpretation}

The operator $M$ defines a canonical element of the Bose--Mesner algebra
of $G$, with
\[
M M^{\mathsf T} = p(A).
\]

Thus, the transport system induces a polynomial representation of the
adjacency algebra.

\subsection{The Cubic Thalean Representation}

\begin{definition}
The map
\[
A \;\mapsto\; A^3 + 2A^2 + 2I
\]
is called the \emph{cubic Thalean representation}.
\end{definition}

This representation is canonical: it is determined entirely by the
transport-induced cocycle and the combinatorial structure of $G$.

\subsection{Spectral Interpretation}

If $Av = \lambda v$, then
\[
M M^{\mathsf T} v = p(\lambda)v.
\]

Thus, the representation acts by spectral transformation via $p(x)$.
